\chapter{Algebraic Varieties}\label{varieties}
\chapterstatus{complete}

\section{Introduction}\label{varieties:section-introduction}

This chapter develops classical algebraic geometry over an algebraically closed field: affine and projective varieties, their regular
functions and morphisms, dimension, tangent spaces and smoothness, rational maps, and blow-ups. The Nullstellensatz of
Chapter~\ref{commutative-algebra} translates geometry into commutative algebra. Closed subsets correspond to radical ideals, varieties to
prime ideals, points to maximal ideals, and affine varieties to finitely generated domains. Chapter~\ref{schemes} generalises everything to
schemes, which are needed for non-closed fields, for families, and for nilpotents. The classical language is enough for most statements
about complex projective manifolds, and it is the one in which the Hodge conjecture is usually stated. References:
\cite[Chapter I]{hartshorne-ag}, \cite[Chapter 1]{griffiths-harris}.

\noindent\textbf{Prerequisites.} Chapter~\ref{commutative-algebra} (Commutative Algebra), Chapter~\ref{general-topology} (General Topology)
and Chapter~\ref{sheaves} (Sheaves on Topological Spaces).

Throughout, $k$ is an algebraically closed field, of any characteristic. Rings are commutative with $1$, and $k[x_1, \ldots, x_n]$ is the
polynomial ring, with its usual addition and multiplication.

\section{Affine algebraic sets and the Zariski topology}\label{varieties:section-affine}

\begin{definition}\label{varieties:definition-zariski}
\emph{Affine $n$-space} is $\Aa^n = k^n$. For a subset $T \subseteq A = k[x_1, \ldots, x_n]$ put $Z(T) = \{a \in \Aa^n : f(a) = 0 \text{ for all } f \in T\}$,
and for $Y \subseteq \Aa^n$ put $I(Y) = \{f \in A : f(a) = 0 \text{ for all } a \in Y\}$. The sets $Z(T)$ are the \emph{algebraic sets}.
\end{definition}

\begin{proposition}\label{varieties:proposition-zariski-topology}
\begin{enumerate}
\item $Z(0) = \Aa^n$, $Z(1) = \emptyset$, $Z(T) = Z((T))$, and $Z(I) \cup Z(J) = Z(IJ) = Z(I \cap J)$ and $\bigcap_\alpha Z(I_\alpha) = Z(\sum_\alpha I_\alpha)$ for ideals.
So the algebraic sets are the closed sets of a topology on $\Aa^n$, the \emph{Zariski topology}.
\item $I(Y)$ is a radical ideal, $Z(I(Y)) = \overline Y$, and $I(Z(J)) = \sqrt J$ for every ideal $J$.
\item $Y \mapsto I(Y)$ and $J \mapsto Z(J)$ are mutually inverse inclusion-reversing bijections between closed subsets of $\Aa^n$ and radical
ideals of $A$. Under them, points correspond to maximal ideals.
\end{enumerate}
\end{proposition}

\begin{proof}
(1) If $a \notin Z(I) \cup Z(J)$, pick $f \in I$, $g \in J$ with $f(a) \neq 0 \neq g(a)$; then $fg \in IJ$ does not vanish at $a$. So
$Z(IJ) \subseteq Z(I) \cup Z(J) \subseteq Z(I \cap J) \subseteq Z(IJ)$. The rest is immediate. (2) If $f^m$ vanishes on $Y$, so does $f$, so $I(Y)$ is
radical. $Z(I(Y))$ is a closed set containing $Y$, and if $Z(J) \supseteq Y$ then $J \subseteq I(Y)$, so $Z(J) \supseteq Z(I(Y))$; hence
$Z(I(Y)) = \overline Y$. The equality $I(Z(J)) = \sqrt J$ is the Nullstellensatz (Theorem~\ref{commutative-algebra:theorem-nullstellensatz}(3)),
the inclusion $\supseteq$ being clear. (3) follows from (2), and points correspond to maximal ideals by
Theorem~\ref{commutative-algebra:theorem-nullstellensatz}(1).
\end{proof}

\begin{definition}\label{varieties:definition-irreducible}
A topological space $Y$ is \emph{irreducible} if it is nonempty and not the union of two proper closed subsets; equivalently, every
nonempty open subset is dense. It is \emph{Noetherian} if every descending chain of closed subsets is eventually constant.
\end{definition}

\begin{proposition}\label{varieties:proposition-irreducible-components}
\begin{enumerate}
\item A closed set $Y \subseteq \Aa^n$ is irreducible if and only if $I(Y)$ is prime.
\item $\Aa^n$, and every subspace of it, is Noetherian.
\item Every Noetherian space $Y$ is a finite union $Y = Y_1 \cup \cdots \cup Y_r$ of irreducible closed subsets, unique up to order if
$Y_i \not\subseteq Y_j$ for $i \neq j$; the $Y_i$ are the \emph{irreducible components}.
\end{enumerate}
\end{proposition}

\begin{proof}
(1) If $Y = Y_1 \cup Y_2$ with $Y_i$ proper closed, pick $f_i \in I(Y_i) \setminus I(Y)$ by Proposition~\ref{varieties:proposition-zariski-topology}(3).
Then $f_1f_2 \in I(Y)$, so $I(Y)$ is not prime. Conversely, if $fg \in I(Y)$, then $Y = (Y \cap Z(f)) \cup (Y \cap Z(g))$, so $Y \subseteq Z(f)$ or
$Y \subseteq Z(g)$ by irreducibility. (Also $Y \neq \emptyset$ if and only if $I(Y) \neq A$.) (2) A descending chain of closed sets gives an
ascending chain of ideals, which stabilises because $A$ is Noetherian (Theorem~\ref{commutative-algebra:theorem-hilbert-basis}). For a subspace,
its closed sets are traces of closed sets of $\Aa^n$. (3) If some closed subset is not such a finite union, choose a minimal one $Y'$, which
exists by the Noetherian property. It is not irreducible, so $Y' = Y'_1 \cup Y'_2$ with smaller closed sets, which are finite unions, a
contradiction. For uniqueness, if $Y_1 \cup \cdots \cup Y_r = Y'_1 \cup \cdots \cup Y'_s$ are irredundant, then
$Y'_1 = \bigcup_i(Y'_1 \cap Y_i)$ forces $Y'_1 \subseteq Y_i$ for some $i$, and similarly $Y_i \subseteq Y'_j$, so $j = 1$ and $Y'_1 = Y_i$; induct.
\end{proof}

\begin{definition}\label{varieties:definition-affine-variety}
An \emph{affine variety} is an irreducible closed subset $Y \subseteq \Aa^n$ with the induced topology. Its \emph{coordinate ring} is
$A(Y) = k[x_1, \ldots, x_n]/I(Y)$, a finitely generated $k$-algebra that is a domain. A \emph{quasi-affine variety} is a nonempty open subset
of an affine variety.
\end{definition}

\begin{example}\label{varieties:example-affine}
\begin{enumerate}
\item $\Aa^n$ is an affine variety, since $I(\Aa^n) = 0$ is prime ($k$ is infinite, so a polynomial vanishing everywhere is $0$).
\item If $f$ is irreducible, then $(f)$ is prime because $k[x_1, \ldots, x_n]$ is a unique factorisation domain
(Theorem~\ref{rings:theorem-polynomial-ufd}), so the \emph{hypersurface} $Z(f)$ is an affine variety with $I(Z(f)) = (f)$.
\item The closed subsets of $\Aa^1$ are $\Aa^1$ and the finite sets, since $k[x]$ is a principal ideal domain. So the Zariski topology is not
Hausdorff: any two nonempty open sets meet.
\item The Zariski topology on $\Aa^2$ is not the product topology on $\Aa^1 \times \Aa^1$. The closed sets of the product topology are the finite
unions of points, horizontal and vertical lines and $\Aa^2$, whereas the parabola $Z(y - x^2)$ is also Zariski closed.
\item The \emph{twisted cubic} $Y = \{(t, t^2, t^3)\} = Z(y - x^2, z - x^3) \subseteq \Aa^3$ is an affine variety with
$A(Y) \cong k[x, y, z]/(y - x^2, z - x^3) \cong k[t]$, a domain.
\end{enumerate}
\end{example}

\section{Regular functions and morphisms}\label{varieties:section-regular}

\begin{definition}\label{varieties:definition-regular-affine}
Let $Y$ be a quasi-affine variety in $\Aa^n$. A function $f \colon Y \to k$ is \emph{regular at} $P \in Y$ if there are an open neighbourhood
$U \ni P$ and polynomials $g, h$ with $h$ nowhere zero on $U$ and $f = g/h$ on $U$. It is \emph{regular} on an open set if it is regular at
each of its points. The regular functions on open subsets form a sheaf of $k$-algebras $\mathcal{O}_Y$. The \emph{local ring} at $P$ is the stalk
$\mathcal{O}_{Y,P}$, and the \emph{function field} $K(Y)$ is the set of regular functions on nonempty open subsets modulo agreement on a
nonempty open subset.
\end{definition}

\begin{lemma}\label{varieties:lemma-regular-continuous}
A regular function $f \colon Y \to k = \Aa^1$ is continuous for the Zariski topologies. If two regular functions on an irreducible $Y$ agree on
a nonempty open subset, they are equal.
\end{lemma}

\begin{proof}
The closed subsets of $\Aa^1$ are finite sets, so it suffices that $f^{-1}(a)$ is closed. Locally $f = g/h$, and
$f^{-1}(a) \cap U = Z(g - ah) \cap U$ is closed in $U$; being closed is a local property. For the second statement, $\{f = f'\} = (f - f')^{-1}(0)$
is closed and contains a dense open subset.
\end{proof}

\begin{theorem}\label{varieties:theorem-regular-functions-affine}
Let $Y$ be an affine variety with coordinate ring $A = A(Y)$. For $P \in Y$ let $\mathfrak{m}_P \subseteq A$ be the maximal ideal of functions
vanishing at $P$. Then:
\begin{enumerate}
\item $\mathcal{O}_{Y,P} \cong A_{\mathfrak{m}_P}$;
\item for $f \in A$, the \emph{basic open set} $D(f) = Y \setminus Z(f)$ satisfies $\mathcal{O}(D(f)) \cong A_f$; in particular
$\mathcal{O}(Y) \cong A$;
\item $K(Y) \cong \mathrm{Frac}(A)$, and the basic open sets form a basis of the topology.
\end{enumerate}
\end{theorem}

\begin{proof}
The basic open sets form a basis because $Y \setminus Z(J) = \bigcup_{f \in J}D(f)$. A polynomial vanishing on $Y$ lies in $I(Y)$, so $A$ is a ring of
functions on $Y$. (1) Send $g/h \in A_{\mathfrak{m}_P}$, $h(P) \neq 0$, to the germ of the function $g/h$ on $D(h)$. This is surjective by definition. It
is injective: if $g/h$ vanishes near $P$, then $g$ vanishes on a nonempty open subset of $Y$, hence on $Y$ by
Lemma~\ref{varieties:lemma-regular-continuous}, so $g = 0$. So all local rings embed compatibly in $\mathrm{Frac}(A)$, and a regular function on
an open set $U$ is determined by its germs, so $\mathcal{O}(U) \subseteq \bigcap_{P \in U}A_{\mathfrak{m}_P} \subseteq \mathrm{Frac}(A)$, with
equality by the definition of regularity.

(2) The points of $D(f)$ correspond to the maximal ideals of $A$ not containing $f$, which are the maximal ideals of $A_f$
(Proposition~\ref{commutative-algebra:proposition-localisation}). So we must show $\bigcap_{\mathfrak{m} \not\ni f}A_\mathfrak{m} = A_f$ inside
$\mathrm{Frac}(A)$. The inclusion $\supseteq$ is clear. Let $\varphi$ lie in the intersection, and let $J = \{a \in A_f : a\varphi \in A_f\}$, an
ideal of $A_f$. For each maximal ideal $\mathfrak{n}$ of $A_f$ we have $\varphi = g/h$ with $h \notin \mathfrak{n}$, so $h \in J \setminus \mathfrak{n}$. Thus
$J$ lies in no maximal ideal, $1 \in J$, and $\varphi \in A_f$. Taking $f = 1$ gives $\mathcal{O}(Y) = A$.

(3) Every nonempty open set contains a nonempty basic open set, and $\mathcal{O}(D(f)) = A_f$ with restriction maps the inclusions of
subrings of $\mathrm{Frac}(A)$. So $K(Y) = \bigcup_{f \neq 0}A_f = \mathrm{Frac}(A)$.
\end{proof}

\begin{definition}\label{varieties:definition-morphism}
Let $X$, $Y$ be varieties, meaning for now quasi-affine varieties (Definition~\ref{varieties:definition-quasi-projective} below extends
this to quasi-projective ones). A \emph{morphism} $\varphi \colon X \to Y$ is a continuous map such that for every open $V \subseteq Y$ and every
regular $f$ on $V$, the function $f \circ \varphi$ is regular on $\varphi^{-1}(V)$. Varieties and morphisms form a category, and an isomorphism is
a morphism with an inverse morphism. Equivalently, a morphism is a morphism of ringed spaces $(X, \mathcal{O}_X) \to (Y, \mathcal{O}_Y)$ whose
map of sheaves is composition of functions (Definition~\ref{sheaves:definition-ringed-space}).
\end{definition}

\begin{proposition}\label{varieties:proposition-morphisms-to-affine}
Let $X$ be a variety and $Y \subseteq \Aa^n$ an affine variety. Then $\varphi \mapsto \varphi^*$, $\varphi^*(f) = f \circ \varphi$, is a bijection
\[
\Hom(X, Y) \to \Hom_{k\text{-alg}}(A(Y), \mathcal{O}(X)).
\]
Consequently $Y \mapsto A(Y)$ is an anti-equivalence between the category of affine varieties and that of finitely generated
$k$-algebras that are domains.
\end{proposition}

\begin{proof}
Given a $k$-algebra map $\theta \colon A(Y) \to \mathcal{O}(X)$, let $\xi_i = \theta(\bar x_i)$ and define $\psi(P) = (\xi_1(P), \ldots, \xi_n(P))$.
For $f \in I(Y)$, $f(\psi(P)) = \theta(\bar f)(P) = 0$, so $\psi$ maps into $Y$, and $f \circ \psi = \theta(\bar f)$ for all polynomials $f$. Then
$\psi^{-1}(Y \cap Z(f)) = \theta(\bar f)^{-1}(0)$ is closed (Lemma~\ref{varieties:lemma-regular-continuous}), so $\psi$ is continuous. A function
that is locally $g/h$ pulls back to $(g \circ \psi)/(h \circ \psi)$, a quotient of regular functions, hence regular, so $\psi$ is a morphism with
$\psi^* = \theta$. Conversely, a morphism $\varphi$ is determined by the functions $x_i \circ \varphi = \varphi^*(\bar x_i)$. For the
anti-equivalence, use $\mathcal{O}(X) = A(X)$ for affine $X$ (Theorem~\ref{varieties:theorem-regular-functions-affine}). Every finitely generated
domain $B = k[x_1, \ldots, x_n]/\mathfrak{p}$ is $A(Z(\mathfrak{p}))$, since $I(Z(\mathfrak{p})) = \sqrt{\mathfrak{p}} = \mathfrak{p}$.
\end{proof}

\begin{example}\label{varieties:example-morphisms}
\begin{enumerate}
\item $\Aa^1 \setminus \{0\}$ is isomorphic to the affine variety $Z(xy - 1) \subseteq \Aa^2$, with coordinate ring $k[x, x^{-1}]$. More generally
$D(f) \subseteq Y$ is isomorphic to the affine variety with coordinate ring $A(Y)_f \cong A(Y)[t]/(tf - 1)$.
\item $\Aa^2 \setminus \{0\}$ is not isomorphic to an affine variety. Its regular functions are
$\mathcal{O}(D(x)) \cap \mathcal{O}(D(y)) = k[x, y]_x \cap k[x, y]_y = k[x, y]$ by Theorem~\ref{varieties:theorem-regular-functions-affine}. So if it were
affine, the inclusion into $\Aa^2$ would induce an isomorphism of coordinate rings and would therefore be an isomorphism by
Proposition~\ref{varieties:proposition-morphisms-to-affine}, which it is not.
\item The morphism $\Aa^1 \to Z(y^2 - x^3)$, $t \mapsto (t^2, t^3)$, is a bijective morphism that is not an isomorphism: it corresponds to the
inclusion $k[t^2, t^3] \subseteq k[t]$, which is not surjective.
\end{enumerate}
\end{example}

\section{Projective varieties}\label{varieties:section-projective}

\begin{definition}\label{varieties:definition-projective-space}
\emph{Projective $n$-space} is $\PP^n = (k^{n+1} \setminus \{0\})/k^\times$, with points written $[a_0 : \cdots : a_n]$. Let $S = k[x_0, \ldots, x_n]$
with its grading by degree, and $S_+ = (x_0, \ldots, x_n)$. For a set $T$ of homogeneous polynomials, $Z(T) \subseteq \PP^n$ is the set of points at
which all $f \in T$ vanish (this is well defined since $f(\lambda a) = \lambda^{\deg f}f(a)$). For $Y \subseteq \PP^n$, $I(Y)$ is the ideal generated
by the homogeneous polynomials vanishing on $Y$, a homogeneous ideal. The sets $Z(\mathfrak{a})$, $\mathfrak{a}$ a homogeneous ideal, are the closed
sets of the \emph{Zariski topology} on $\PP^n$, as in Proposition~\ref{varieties:proposition-zariski-topology}(1).
\end{definition}

\begin{proposition}[Projective Nullstellensatz]\label{varieties:proposition-projective-nullstellensatz}
Let $\mathfrak{a} \subseteq S$ be a homogeneous ideal. Then $Z(\mathfrak{a}) = \emptyset$ if and only if $\sqrt{\mathfrak{a}} \supseteq S_+$, and
otherwise $I(Z(\mathfrak{a})) = \sqrt{\mathfrak{a}}$. Hence closed subsets of $\PP^n$ correspond bijectively to homogeneous radical ideals not
containing $S_+$, and irreducible closed subsets to homogeneous primes not containing $S_+$.
\end{proposition}

\begin{proof}
Let $C \subseteq \Aa^{n+1}$ be the affine zero set of $\mathfrak{a}$. Since $\mathfrak{a}$ is homogeneous, $C$ is a union of lines through $0$ (or
empty), and $Z(\mathfrak{a}) = (C \setminus \{0\})/k^\times$. So $Z(\mathfrak{a}) = \emptyset$ if and only if $C \subseteq \{0\}$, if and only if
$\sqrt{\mathfrak{a}} \supseteq I(\{0\}) = S_+$ by the affine Nullstellensatz. Otherwise $C \neq \{0\}$ contains $0$ if $\mathfrak{a} \neq S$, and a
homogeneous $f$ of positive degree vanishes on $Z(\mathfrak{a})$ if and only if it vanishes on $C$, if and only if $f \in \sqrt{\mathfrak{a}}$. The
radical of a homogeneous ideal is homogeneous: if $f^m \in \mathfrak{a}$ and $f_d$ is the top component of $f$, then $f_d^m$ is the top component of
$f^m$, so it lies in $\mathfrak{a}$ (Lemma~\ref{rings:lemma-homogeneous-ideal}), and $f_d \in \sqrt{\mathfrak{a}}$; now induct on the number of
components of $f - f_d$. Hence $I(Z(\mathfrak{a})) = \sqrt{\mathfrak{a}}$.

For irreducibility, argue as in Proposition~\ref{varieties:proposition-irreducible-components}(1) with homogeneous $f$, $g$. This suffices
because a homogeneous ideal $\mathfrak{p}$ such that $fg \in \mathfrak{p}$ implies $f \in \mathfrak{p}$ or $g \in \mathfrak{p}$ for homogeneous $f$, $g$
is prime. Indeed, let $fg \in \mathfrak{p}$ with $f, g \notin \mathfrak{p}$. Subtracting from $f$ and $g$ their homogeneous components of high
degree that lie in $\mathfrak{p}$ does not change the hypotheses, so we may assume that the top components $f_d$ and $g_e$ are not in
$\mathfrak{p}$. But $f_dg_e$ is the top component of $fg$, so it lies in $\mathfrak{p}$, a contradiction.
\end{proof}

\begin{proposition}\label{varieties:proposition-standard-cover}
Let $U_i = \PP^n \setminus Z(x_i)$. The map $\varphi_i \colon U_i \to \Aa^n$, $[a] \mapsto (a_j/a_i)_{j \neq i}$, is a homeomorphism for the Zariski
topologies. Under it, closed subsets of $U_i$ correspond to closed subsets of $\Aa^n$ through dehomogenisation
$f \mapsto f(x_0, \ldots, 1, \ldots, x_n)$ and homogenisation.
\end{proposition}

\begin{proof}
$\varphi_i$ is bijective with inverse $(b_j) \mapsto [b_0 : \cdots : 1 : \cdots : b_n]$. For homogeneous $f$ of degree $d$, $\varphi_i(Z(f) \cap U_i)$ is
the zero set of the dehomogenisation $f(x_0, \ldots, 1, \ldots, x_n)$. Conversely, for $g \in k[y_j : j \neq i]$ of degree $d$, the homogenisation
$x_i^dg(x_0/x_i, \ldots)$ is a homogeneous polynomial whose zero set in $U_i$ is $\varphi_i^{-1}(Z(g))$. So $\varphi_i$ and its inverse are closed maps.
\end{proof}

\begin{definition}\label{varieties:definition-quasi-projective}
A \emph{projective variety} is an irreducible closed subset $Y \subseteq \PP^n$, and its \emph{homogeneous coordinate ring} is
$S(Y) = S/I(Y)$, a graded domain. A \emph{quasi-projective variety} is a nonempty open subset of a projective variety. A function
$f \colon Y \to k$ on a quasi-projective variety is \emph{regular at} $P$ if near $P$ it is $g/h$ with $g$, $h$ homogeneous of the same degree and $h$
nowhere zero. This gives the sheaf $\mathcal{O}_Y$, the local rings $\mathcal{O}_{Y,P}$ and the function field $K(Y)$ as before. From now on a
\emph{variety} is a quasi-projective variety, and morphisms are defined as in Definition~\ref{varieties:definition-morphism}.
\end{definition}

\begin{proposition}\label{varieties:proposition-affine-cover}
Let $Y \subseteq \PP^n$ be a quasi-projective variety. Then $\varphi_i$ restricts to an isomorphism of $Y \cap U_i$ (if nonempty) onto a
quasi-affine variety, and onto an affine variety if $Y$ is projective, with coordinate ring the degree-zero part $S(Y)_{(x_i)}$ of $S(Y)_{x_i}$.
In particular, every quasi-affine variety is quasi-projective (via $\Aa^n \cong U_0$), every variety has an open cover by affine varieties, and
$K(Y) \cong K(Y \cap U_i)$.
\end{proposition}

\begin{proof}
By Proposition~\ref{varieties:proposition-standard-cover}, $\varphi_i$ is a homeomorphism onto its image, which is locally closed, and irreducible
since $Y \cap U_i$ is a nonempty open subset of an irreducible space. A quotient $g/h$ of homogeneous polynomials of degree $d$ equals the
quotient of the dehomogenisations of $g$ and $h$ on $U_i$, and conversely a quotient of polynomials in the affine coordinates can be homogenised
to a quotient of homogeneous polynomials of the same degree. So regular functions correspond. If $Y$ is projective, $\varphi_i(Y \cap U_i)$ is
closed, and dehomogenisation identifies its ideal with the degree-zero part of the localised ideal. A quasi-affine variety is covered by basic
open sets $D(f)$ of its closure, which are affine (Example~\ref{varieties:example-morphisms}(1)). Finally, every nonempty open subset of $Y$ meets
$U_i$ in a dense open subset, so the function fields agree.
\end{proof}

\begin{theorem}\label{varieties:theorem-projective-constant}
If $Y$ is a projective variety, then $\mathcal{O}(Y) = k$. Consequently every morphism from a projective variety to an affine variety is
constant.
\end{theorem}

\begin{proof}
Let $f \in \mathcal{O}(Y)$ and $L = \mathrm{Frac}(S(Y))$. We view $K(Y)$ as the degree-zero part of $L$, by $g/h \mapsto g/h$. Let $\Sigma$ be the set of $i$ with $x_i \notin I(Y)$;
the other $x_i$ are zero in $S(Y)$. For $i \in \Sigma$, $f$ restricted to $Y \cap U_i$ lies in $S(Y)_{(x_i)}$ by
Proposition~\ref{varieties:proposition-affine-cover} and Theorem~\ref{varieties:theorem-regular-functions-affine}, so $x_i^{N_i}f \in S(Y)_{N_i}$ for
some $N_i$, as elements of $L$. Let $N = \sum_{i \in \Sigma}N_i$. The space $S(Y)_N$ is spanned by the monomials of degree $N$ in the $x_i$, $i \in \Sigma$,
and each of them is divisible by $x_i^{N_i}$ for some $i \in \Sigma$. So $S(Y)_N \cdot f \subseteq S(Y)_N$, and by induction
$S(Y)_N \cdot f^q \subseteq S(Y)_N$ for all $q \geq 1$. Choose $i \in \Sigma$. Then
$x_i^Nf^q \in S(Y)$ for all $q$, so $S(Y)[f] \subseteq x_i^{-N}S(Y)$, a finitely generated $S(Y)$-module. As $S(Y)$ is Noetherian, $S(Y)[f]$ is
a finitely generated $S(Y)$-module, and $f$ is integral over $S(Y)$ (Proposition~\ref{commutative-algebra:proposition-integral}(1)). Taking the
degree-zero parts of the coefficients of an integral equation shows that $f$ is integral over $S(Y)_0 = k$. As $k$ is algebraically closed and
$f$ lies in a field, $f \in k$. For the second statement, compose with the coordinate functions of the target.
\end{proof}

\begin{example}\label{varieties:example-projective}
\begin{enumerate}
\item A \emph{projective hypersurface} is $Z(f) \subseteq \PP^n$ with $f$ irreducible and homogeneous. Plane curves are the case $n = 2$.
\item The \emph{$d$-uple (Veronese) embedding} $v_d \colon \PP^n \to \PP^N$, $N = \binom{n+d}{n} - 1$, sends $a$ to the vector of all monomials of degree
$d$ in $a$. Its image is closed, cut out by the quadratic relations $z_Mz_{M'} = z_{M''}z_{M'''}$ whenever $MM' = M''M'''$, and $v_d$ is an isomorphism
onto it \cite[Chapter I, Exercises 2.12 and 3.4]{hartshorne-ag}. On the image, the inverse is given on $\{z_{x_i^d} \neq 0\}$ by
$z \mapsto [z_{x_i^{d-1}x_0} : \cdots : z_{x_i^{d-1}x_n}]$. For $n = 1$, $d = 3$ the image is the \emph{twisted cubic} in $\PP^3$, the closure of the affine twisted cubic.
\item The \emph{Segre embedding} $\PP^m \times \PP^n \to \PP^{mn+m+n}$, $(a, b) \mapsto (a_ib_j)$, has closed image $Z(z_{ij}z_{kl} - z_{il}z_{kj})$. This image
is used to give $\PP^m \times \PP^n$ the structure of a projective variety. For $m = n = 1$ it is the smooth quadric surface $Z(z_{00}z_{11} - z_{01}z_{10})$
in $\PP^3$.
\end{enumerate}
\end{example}

\begin{proposition}\label{varieties:proposition-products-affine}
Let $X \subseteq \Aa^m$ and $Y \subseteq \Aa^n$ be affine varieties. Then $X \times Y \subseteq \Aa^{m+n}$ is an affine variety with coordinate ring
$A(X) \otimes_k A(Y)$, and it is the product of $X$ and $Y$ in the category of varieties.
\end{proposition}

\begin{proof}
$X \times Y = Z(I(X) + I(Y))$ is closed. It is irreducible: if $X \times Y = Z_1 \cup Z_2$ with $Z_i$ closed, each fibre $\{x\} \times Y \cong Y$ is
irreducible, so it lies in $Z_1$ or in $Z_2$. The set $X_i = \{x : \{x\} \times Y \subseteq Z_i\} = \bigcap_{y \in Y}\{x : (x, y) \in Z_i\}$ is closed,
and $X = X_1 \cup X_2$, so $X = X_1$ or $X = X_2$. The natural map $A(X) \otimes_k A(Y) \to \{\text{functions on } X \times Y\}$ is injective:
write an element as $\sum_ia_i \otimes b_i$ with the $b_i$ linearly independent over $k$; if it vanishes, then for each $x$ the function
$\sum_ia_i(x)b_i$ vanishes on $Y$, so all $a_i(x) = 0$, so all $a_i = 0$. Its image is the ring of polynomial functions on $X \times Y$, which is
$A(X \times Y)$. The universal property follows from Proposition~\ref{varieties:proposition-morphisms-to-affine} and the universal property of
$\otimes_k$ for $k$-algebras (Proposition~\ref{rings:proposition-tensor-algebras}).
\end{proof}

\section{Dimension}\label{varieties:section-dimension}

\begin{definition}\label{varieties:definition-dimension}
The \emph{dimension} of a topological space is the supremum of the lengths $r$ of chains $Z_0 \subsetneq Z_1 \subsetneq \cdots \subsetneq Z_r$ of
irreducible closed subsets. For an irreducible closed subset $Z \subseteq Y$, the \emph{codimension} of $Z$ in $Y$ is the supremum of the lengths of
such chains with $Z_0 = Z$. Varieties of dimension $1$ and $2$ are \emph{curves} and \emph{surfaces}.
\end{definition}

\begin{lemma}\label{varieties:lemma-catenary}
Let $A$ be a finitely generated $k$-algebra that is a domain, of dimension $d$. Then every maximal chain of primes of $A$ has length $d$.
Consequently $\operatorname{ht}\mathfrak{p} + \dim A/\mathfrak{p} = d$ for every prime $\mathfrak{p}$, and every maximal ideal has height $d$.
\end{lemma}

\begin{proof}
Induction on $d = \operatorname{trdeg}_k\mathrm{Frac}(A)$ (Theorem~\ref{commutative-algebra:theorem-dimension-finitely-generated}). If $d = 0$,
$A$ is a domain algebraic over $k$, hence a field (Proposition~\ref{commutative-algebra:proposition-integral}(3)), and the only chain is $0$. Let
$d > 0$ and let $0 = \mathfrak{p}_0 \subsetneq \mathfrak{p}_1 \subsetneq \cdots \subsetneq \mathfrak{p}_m$ be a maximal chain. Then $m \geq 1$, and no prime
lies strictly between $0$ and $\mathfrak{p}_1$. By Noether normalisation (Theorem~\ref{commutative-algebra:theorem-noether-normalisation}), $A$ is
finite over a polynomial ring $B = k[y_1, \ldots, y_d]$. Let $\mathfrak{q} = \mathfrak{p}_1 \cap B$.

The prime $\mathfrak{q}$ is nonzero. Indeed, take $a \in \mathfrak{p}_1 \setminus 0$ and an equation $a^r + b_1a^{r-1} + \cdots + b_r = 0$ over $B$ with $r$
minimal. Then $b_r \neq 0$, since otherwise we could divide by $a$, and $b_r \in aA \cap B \subseteq \mathfrak{q}$. If there were a prime
$0 \subsetneq \mathfrak{q}' \subsetneq \mathfrak{q}$, then going down for the integral extension $B \subseteq A$ of domains with $B$ normal
(Remark~\ref{commutative-algebra:remark-going-down}; $B$ is a unique factorisation domain, hence normal by
Proposition~\ref{commutative-algebra:proposition-integral}(4)) would give a prime $\mathfrak{p}' \subsetneq \mathfrak{p}_1$ lying over $\mathfrak{q}'$, so
$\mathfrak{p}' \neq 0$, which is a contradiction. So $\mathfrak{q}$ is a minimal nonzero prime of the unique factorisation domain $B$. It contains an
irreducible factor $f$ of any of its nonzero elements, and $(f)$ is a nonzero prime, so $\mathfrak{q} = (f)$. If $f$ involves $y_i$, the images of the
other $y_j$ in $B/(f)$ are algebraically independent, since a relation $g(y_j : j \neq i) \in (f)$ is impossible for $g \neq 0$, and $y_i$ is
algebraic over them. So $\operatorname{trdeg}\mathrm{Frac}(B/\mathfrak{q}) = d - 1$. As $A/\mathfrak{p}_1$ is integral over $B/\mathfrak{q}$, it also has
transcendence degree $d - 1$, and by induction the maximal chain $\mathfrak{p}_1/\mathfrak{p}_1 \subsetneq \cdots \subsetneq \mathfrak{p}_m/\mathfrak{p}_1$ of
$A/\mathfrak{p}_1$ has length $d - 1$. Hence $m = d$.

For a prime $\mathfrak{p}$, a chain of length $\operatorname{ht}\mathfrak{p}$ ending at $\mathfrak{p}$ and a chain of length $\dim A/\mathfrak{p}$ starting at
$\mathfrak{p}$ cannot be refined, since they have maximal length. Their concatenation is therefore a maximal chain of $A$, of length $d$.
\end{proof}

\begin{theorem}\label{varieties:theorem-dimension}
Let $Y$ be a variety.
\begin{enumerate}
\item $\dim Y = \operatorname{trdeg}_kK(Y)$; if $Y$ is affine, $\dim Y = \dim A(Y)$.
\item $\dim U = \dim Y$ for every nonempty open $U \subseteq Y$, and $\dim\mathcal{O}_{Y,P} = \dim Y$ for every $P \in Y$.
\item If $Z \subsetneq Y$ is closed and irreducible, then $\dim Z < \dim Y$, and $\dim Z + \operatorname{codim}(Z, Y) = \dim Y$.
\end{enumerate}
\end{theorem}

\begin{proof}
First let $Y$ be affine with $A = A(Y)$. Irreducible closed subsets of $Y$ correspond to primes of $A$, reversing inclusions
(Propositions~\ref{varieties:proposition-zariski-topology} and~\ref{varieties:proposition-irreducible-components}), so
$\dim Y = \dim A = \operatorname{trdeg}\mathrm{Frac}(A) = \operatorname{trdeg}K(Y)$ by
Theorem~\ref{commutative-algebra:theorem-dimension-finitely-generated} and Theorem~\ref{varieties:theorem-regular-functions-affine}. Also
$\dim\mathcal{O}_{Y,P} = \dim A_{\mathfrak{m}_P} = \operatorname{ht}\mathfrak{m}_P = \dim Y$ by Lemma~\ref{varieties:lemma-catenary}. For $U$ open,
$Z \mapsto \overline Z$ maps chains of irreducible closed subsets of $U$ injectively to chains in $Y$, so $\dim U \leq \dim Y$. Conversely, let
$P \in U$ and take a chain of primes of length $\dim Y$ contained in $\mathfrak{m}_P$. It gives a chain of irreducible closed subsets of $Y$ all
containing $P$, and their intersections with $U$ form a chain of the same length in $U$, since a nonempty open subset of an irreducible space is
dense. Statement (3) for affine $Y$ is Lemma~\ref{varieties:lemma-catenary}, since $\operatorname{codim}(Z, Y) = \operatorname{ht}I(Z)/I(Y)$ and
$\operatorname{ht} \geq 1$ for a nonzero prime.

In general, cover $Y$ by affine open subsets $Y_i$ (Proposition~\ref{varieties:proposition-affine-cover}); they have
$K(Y_i) = K(Y)$, so $\dim Y_i = \operatorname{trdeg}K(Y)$. Given a chain $Z_0 \subsetneq \cdots \subsetneq Z_r$ in $Y$, choose $Y_i$ meeting $Z_0$; the
sets $Z_j \cap Y_i$ form a chain in $Y_i$ (their closures are the $Z_j$). So $\dim Y = \max_i\dim Y_i = \operatorname{trdeg}K(Y)$, and the other
statements reduce to the affine case in the same way: a proper closed irreducible $Z$ meets some $Y_i$ in a proper subset, since otherwise $Z$
would contain the dense set $Y_i$.
\end{proof}

\begin{theorem}[Krull's theorem for varieties]\label{varieties:theorem-hypersurface-sections}
Let $Y$ be an affine variety of dimension $r$ and $f \in A(Y)$ neither zero nor a unit. Then every irreducible component of $Y \cap Z(f)$ has
dimension $r - 1$. A variety $Y \subseteq \Aa^n$ has dimension $n - 1$ if and only if $Y = Z(f)$ for an irreducible polynomial $f$.
\end{theorem}

\begin{proof}
The components of $Z(f)$ correspond to the minimal primes over $(f)$ in $A(Y)$. They have height $\leq 1$ by Krull's principal ideal theorem
(Theorem~\ref{commutative-algebra:theorem-krull-hauptidealsatz}), hence height $1$ as they are nonzero, hence dimension $r - 1$ by
Lemma~\ref{varieties:lemma-catenary}. For the second statement, $Z(f)$ has dimension $n - 1$ by the first. Conversely, if $\dim Y = n - 1$, then
$I(Y)$ has height $1$. It contains an irreducible factor $f$ of any of its nonzero elements, and $0 \subsetneq (f) \subseteq I(Y)$ with $(f)$ prime,
so $I(Y) = (f)$.
\end{proof}

\begin{theorem}[Projective dimension theorem]\label{varieties:theorem-projective-dimension}
Let $Y \subseteq \PP^n$ be a projective variety of dimension $r \geq 1$ and $f$ a homogeneous polynomial of positive degree with $f \notin I(Y)$.
Then $Y \cap Z(f)$ is nonempty and all its irreducible components have dimension $r - 1$.
\end{theorem}

\begin{proof}
Let $C(Y) \subseteq \Aa^{n+1}$ be the affine cone, the affine variety with coordinate ring $S(Y)$. The fraction field of $S(Y)$ is $K(Y)(x_i)$
for any $x_i \notin I(Y)$, with $x_i$ transcendental over $K(Y)$ (elements of $K(Y)$ have degree zero), so $\dim C(Y) = r + 1$ by
Theorem~\ref{varieties:theorem-dimension}. The same holds for every projective variety, in particular for the components of $Y \cap Z(f)$, whose
cones are the components of $C(Y) \cap Z(f)$ other than $\{0\}$. The set $C(Y) \cap Z(f)$ contains $0$, and $f$ is not a unit in $S(Y)$, so by
Theorem~\ref{varieties:theorem-hypersurface-sections} all its components have dimension $r \geq 1$. In particular they are not $\{0\}$, so they are
cones over components of $Y \cap Z(f)$ of dimension $r - 1$.
\end{proof}

\begin{corollary}\label{varieties:corollary-intersections-nonempty}
Any $r \leq \dim Y$ hypersurfaces in $\PP^n$ meet a projective variety $Y$. In particular two curves in $\PP^2$ always meet, and $\PP^2$ is not
isomorphic to $\PP^1 \times \PP^1$, which contains disjoint curves.
\end{corollary}

\begin{proof}
Apply Theorem~\ref{varieties:theorem-projective-dimension} repeatedly to a component of dimension $\geq \dim Y - j$ after $j$ steps. (If some $f$
vanishes on the current component, the dimension does not drop.) The curves $\PP^1 \times \{0\}$ and $\PP^1 \times \{\infty\}$ are disjoint.
\end{proof}

\section{Tangent spaces and smoothness}\label{varieties:section-smoothness}

\begin{definition}\label{varieties:definition-tangent-space}
Let $P$ be a point of a variety $Y$, with local ring $(\mathcal{O}_{Y,P}, \mathfrak{m}_P)$ and residue field $k$. The \emph{Zariski tangent space}
is $T_PY = (\mathfrak{m}_P/\mathfrak{m}_P^2)^\vee$. The variety is \emph{nonsingular} (or \emph{smooth}) at $P$ if $\mathcal{O}_{Y,P}$ is a regular
local ring (Definition~\ref{commutative-algebra:definition-regular}), that is, if $\dim_kT_PY = \dim Y$, and \emph{singular} at $P$ otherwise.
$Y$ is nonsingular if it is nonsingular at every point.
\end{definition}

\begin{lemma}\label{varieties:lemma-tangent-dimension}
For every Noetherian local ring $(R, \mathfrak{m}, k)$, $\dim R \leq \dim_k\mathfrak{m}/\mathfrak{m}^2$. Hence $\dim T_PY \geq \dim Y$.
\end{lemma}

\begin{proof}
By Nakayama's lemma (Corollary~\ref{commutative-algebra:corollary-nakayama-generators}), $\mathfrak{m}$ is generated by
$e = \dim\mathfrak{m}/\mathfrak{m}^2$ elements, and $\mathfrak{m}$ is minimal over them, so $\dim R = \operatorname{ht}\mathfrak{m} \leq e$ by
Theorem~\ref{commutative-algebra:theorem-krull-hauptidealsatz}. The second statement uses $\dim\mathcal{O}_{Y,P} = \dim Y$.
\end{proof}

\begin{theorem}[Jacobian criterion]\label{varieties:theorem-jacobian}
Let $Y \subseteq \Aa^n$ be an affine variety of dimension $r$ with $I(Y) = (f_1, \ldots, f_t)$, and $P \in Y$. Then
$\dim T_PY = n - \operatorname{rank}J(P)$, where $J = (\partial f_i/\partial x_j)$ is the Jacobian matrix. So $Y$ is nonsingular at $P$ if and only
if $\operatorname{rank}J(P) = n - r$.
\end{theorem}

\begin{proof}
Let $M = (x_1 - P_1, \ldots, x_n - P_n) \subseteq k[x]$. The linear map $\theta \colon M \to k^n$, $g \mapsto (\partial g/\partial x_j(P))_j$, is onto with
kernel $M^2$ (expand $g$ in powers of $x - P$), so $M/M^2 \cong k^n$. The maximal ideal $\mathfrak{m}$ of $A(Y)$ at $P$ satisfies
$\mathfrak{m}/\mathfrak{m}^2 \cong M/(M^2 + I(Y))$, and localising at $\mathfrak{m}$ does not change $\mathfrak{m}/\mathfrak{m}^2$. The image of $I(Y)$ in
$M/M^2 \cong k^n$ is spanned by the $\theta(f_i)$, the rows of $J(P)$, because $\theta(gf_i) = g(P)\theta(f_i)$ as $f_i(P) = 0$. So
$\dim\mathfrak{m}/\mathfrak{m}^2 = n - \operatorname{rank}J(P)$.
\end{proof}

\begin{example}\label{varieties:example-singular}
\begin{enumerate}
\item The cusp $Z(y^2 - x^3)$ and the node $Z(y^2 - x^2(x + 1))$ (characteristic $\neq 2$) are singular exactly at the origin, where the
Jacobian $(\partial f/\partial x, \partial f/\partial y)$ vanishes.
\item A hypersurface $Z(f) \subseteq \Aa^n$ is singular exactly at the common zeros of $f$ and all $\partial f/\partial x_j$.
\item $\Aa^n$, $\PP^n$ and the smooth quadric $Z(xw - yz) \subseteq \PP^3$ are nonsingular.
\end{enumerate}
\end{example}

\section{Rational maps and birational equivalence}\label{varieties:section-rational-maps}

\begin{lemma}\label{varieties:lemma-morphisms-agree}
Two morphisms $\varphi, \psi \colon X \to Y$ of varieties that agree on a nonempty open subset are equal.
\end{lemma}

\begin{proof}
Embed $Y$ in some $\PP^n$. The set $E = \{x : \varphi(x) = \psi(x)\}$ is closed. Indeed, near a point of $X$, $\varphi$ and $\psi$ map into standard open
sets $U_a$ and $U_b$, so on a neighbourhood $V$ they are given by regular functions $\varphi = [f_0 : \cdots : f_n]$ with $f_a = 1$ and
$\psi = [g_0 : \cdots : g_n]$ with $g_b = 1$. On $V$, $E$ is the common zero set of the regular functions $f_ig_j - f_jg_i$. The closed set $E$
contains a dense open subset of the irreducible $X$, so $E = X$.
\end{proof}

\begin{definition}\label{varieties:definition-rational-map}
A \emph{rational map} $\varphi \colon X \dashrightarrow Y$ is an equivalence class of pairs $(U, \varphi_U)$, where $U \subseteq X$ is nonempty open and
$\varphi_U \colon U \to Y$ is a morphism, two pairs being equivalent if the morphisms agree on the intersection (which is possible by
Lemma~\ref{varieties:lemma-morphisms-agree}). It is \emph{dominant} if $\varphi_U(U)$ is dense. A \emph{birational map} is a dominant rational map
with a dominant rational inverse, and $X$, $Y$ are \emph{birational} if there is one. A variety birational to $\PP^n$ is \emph{rational}.
\end{definition}

\begin{theorem}\label{varieties:theorem-birational-fields}
The functor $X \mapsto K(X)$ is an anti-equivalence between the category of varieties with dominant rational maps and the category of
finitely generated field extensions of $k$. In particular $X$ and $Y$ are birational if and only if $K(X) \cong K(Y)$ over $k$, if and only
if they have isomorphic nonempty open subsets.
\end{theorem}

\begin{proof}
A dominant $\varphi$ defined on $U$ pulls back regular functions on nonempty open subsets of $Y$ to regular functions on nonempty open
subsets of $U$, since preimages of dense open sets under a dominant map are nonempty. This gives $\varphi^* \colon K(Y) \to K(X)$. Conversely, let
$\theta \colon K(Y) \to K(X)$ be a $k$-embedding. Choose an affine open $V \subseteq Y$ with $A(V) = k[y_1, \ldots, y_n]$. The elements $\theta(y_i)$
are regular on some nonempty open $U \subseteq X$, and $\theta$ restricts to $A(V) \to \mathcal{O}(U)$. By
Proposition~\ref{varieties:proposition-morphisms-to-affine} this is $\varphi^*$ for a morphism $\varphi \colon U \to V$. It is dominant: if its
image lay in a proper closed subset $Z(g)$ with $0 \neq g \in A(V)$, then $\theta(g) = 0$, contradicting injectivity. These constructions are
inverse to each other by Lemma~\ref{varieties:lemma-morphisms-agree} and Theorem~\ref{varieties:theorem-regular-functions-affine}. Every finitely
generated extension $K = k(y_1, \ldots, y_n)$ is $K(V)$ for the affine variety with coordinate ring $k[y_1, \ldots, y_n] \subseteq K$
(Proposition~\ref{varieties:proposition-morphisms-to-affine}).

If $\varphi \colon X \dashrightarrow Y$ and $\psi \colon Y \dashrightarrow X$ are inverse birational maps, defined on $U$ and $V$, then
$\psi \circ \varphi = \mathrm{id}$ on $\varphi^{-1}(V)$ and $\varphi \circ \psi = \mathrm{id}$ on $\psi^{-1}(U)$. So $\varphi$ restricts to an
isomorphism $\varphi^{-1}(\psi^{-1}(U)) \to \psi^{-1}(U)$. The converse is clear.
\end{proof}

\begin{proposition}\label{varieties:proposition-birational-hypersurface}
Every variety of dimension $r$ is birational to a hypersurface in $\Aa^{r+1}$.
\end{proposition}

\begin{proof}
The finitely generated extension $K(Y)/k$ of transcendence degree $r$ is separably generated: there are $x_1, \ldots, x_r$, algebraically
independent, such that $K(Y)$ is a finite separable extension of $k(x_1, \ldots, x_r)$. In characteristic $0$ this holds for every transcendence
basis. In characteristic $p$ it holds because $k$ is perfect \cite[Section 26]{matsumura-crt}. By the primitive element theorem
(Theorem~\ref{fields:theorem-primitive-element}), $K(Y) = k(x_1, \ldots, x_r, y)$ with $y$ separable algebraic over $k(x)$. Clearing
denominators in the minimal polynomial of $y$ gives an irreducible $f \in k[x_1, \ldots, x_r, t]$ with $f(x, y) = 0$, and $K(Z(f)) \cong K(Y)$.
Apply Theorem~\ref{varieties:theorem-birational-fields}.
\end{proof}

\begin{theorem}\label{varieties:theorem-singular-locus}
The singular points of a variety $Y$ form a proper closed subset.
\end{theorem}

\begin{proof}
Closedness is local, so let $Y \subseteq \Aa^n$ be affine of dimension $r$ with $I(Y) = (f_1, \ldots, f_t)$. By
Lemma~\ref{varieties:lemma-tangent-dimension} and Theorem~\ref{varieties:theorem-jacobian}, $\operatorname{rank}J(P) \leq n - r$ everywhere, and the
singular set is the common zero set of the $(n - r)$-minors of $J$, which is closed. It is proper: by
Proposition~\ref{varieties:proposition-birational-hypersurface} and Theorem~\ref{varieties:theorem-birational-fields}, $Y$ has an open subset
isomorphic to an open subset of a hypersurface $Z(f) \subseteq \Aa^{r+1}$ with $f$ irreducible, so it suffices to find a nonsingular point on $Z(f)$.
The singular points are $Z(f, \partial f/\partial x_1, \ldots)$ (Example~\ref{varieties:example-singular}). If they were all of $Z(f)$, then each
$\partial f/\partial x_j$ would lie in $I(Z(f)) = (f)$, hence vanish, by degree. In characteristic $0$ this makes $f$ constant. In
characteristic $p$ it makes $f$ a polynomial in $x_1^p, \ldots, x_{r+1}^p$, hence a $p$-th power since $k$ is perfect, contradicting
irreducibility. So the nonsingular points of $Z(f)$ form a nonempty open set, which meets the open subset isomorphic to an open subset of $Y$,
because $Z(f)$ is irreducible. Isomorphic open subsets have isomorphic local rings.
\end{proof}

\begin{example}\label{varieties:example-birational}
\begin{enumerate}
\item $\PP^n$ and $\Aa^n$ are birational, and so are $\PP^1 \times \PP^1$ and $\PP^2$: all three have function field $k(x_1, x_2)$ for $n = 2$. By
Corollary~\ref{varieties:corollary-intersections-nonempty}, $\PP^1 \times \PP^1$ and $\PP^2$ are not isomorphic.
\item The cuspidal cubic $Z(y^2 - x^3)$ is rational: $t = y/x$ gives $K = k(t)$. It is singular, so it is not isomorphic to $\PP^1$ or $\Aa^1$.
\item The \emph{Cremona transformation} $\PP^2 \dashrightarrow \PP^2$, $[x : y : z] \mapsto [yz : xz : xy]$, is birational and equal to its own
inverse, but it is not defined at the three coordinate points.
\item A nonsingular plane cubic is not rational. This follows from the genus (Chapter~\ref{curves}): rational nonsingular projective curves are
isomorphic to $\PP^1$, which has genus $0$, while plane cubics have genus $1$.
\end{enumerate}
\end{example}

\section{Blow-ups}\label{varieties:section-blowups}

\begin{definition}\label{varieties:definition-blowup}
The \emph{blow-up} of $\Aa^n$ at $0$ is the closed subset $\mathrm{Bl}_0\Aa^n = \{(x, y) \in \Aa^n \times \PP^{n-1} : x_iy_j = x_jy_i \text{ for all } i, j\}$
with the projection $\pi \colon \mathrm{Bl}_0\Aa^n \to \Aa^n$. For a closed subvariety $Y \subseteq \Aa^n$ through $0$, the \emph{blow-up} (or
\emph{strict transform}) $\tilde Y$ is the closure of $\pi^{-1}(Y \setminus \{0\})$ in $\mathrm{Bl}_0\Aa^n$.
\end{definition}

\begin{proposition}\label{varieties:proposition-blowup}
$\mathrm{Bl}_0\Aa^n$ is a nonsingular variety of dimension $n$. The map $\pi$ is an isomorphism over $\Aa^n \setminus \{0\}$, and
$E = \pi^{-1}(0) = \{0\} \times \PP^{n-1}$ parametrises the lines through $0$. For $Y$ as above, $\tilde Y \to Y$ is birational, and
$\tilde Y \cap E$ consists of limits of secant directions of $Y$ at $0$.
\end{proposition}

\begin{proof}
On $\{y_i \neq 0\}$ put $y_i = 1$. The equations become $x_j = x_iy_j$, so this open set is the graph of a morphism $\Aa^n \to \Aa^{n-1}$, which is
isomorphic to $\Aa^n$ with coordinates $(x_i, y_j)_{j \neq i}$. These charts are irreducible of dimension $n$ and nonsingular, and pairwise they
meet in nonempty open sets, so $\mathrm{Bl}_0\Aa^n$ is irreducible, of dimension $n$ and nonsingular. Over $x \neq 0$, the equations force
$y = [x]$, so $x \mapsto (x, [x])$ is an inverse of $\pi$ there. Over $0$ the equations are empty, so $E = \{0\} \times \PP^{n-1}$. For $Y$, the map
$\tilde Y \to Y$ is an isomorphism over $Y \setminus \{0\}$, which is dense in $\tilde Y$ by definition, and a point $(0, [v]) \in \tilde Y$ is a limit
of points $(x, [x])$ with $x \in Y \setminus \{0\}$.
\end{proof}

\begin{example}\label{varieties:example-blowup-curves}
Let $\mathrm{char}\,k \neq 2$. For the node $Y = Z(y^2 - x^2(x + 1))$, in the chart $y = xt$ the preimage is $Z(x^2(t^2 - x - 1))$, the union of $E = Z(x)$
(with multiplicity $2$) and the strict transform $\tilde Y = Z(t^2 - x - 1)$. The latter is nonsingular and meets $E$ at the two points $t = \pm 1$,
the two tangent directions of the branches: blowing up separates the branches. For the cusp $Z(y^2 - x^3)$, the strict transform is
$Z(t^2 - x)$, nonsingular and meeting $E$ at one point. In both cases one blow-up \emph{resolves} the singularity.
\end{example}

\begin{remark}\label{varieties:remark-resolution}
Blow-ups at points glue to blow-ups of any variety at a nonsingular point, and more generally along closed subvarieties. For $k = \CC$ this
is the construction of Definition~\ref{complex-manifolds:definition-blowup}. Hironaka proved that in characteristic $0$ every variety is
birational to a nonsingular projective variety, and in fact that singularities can be resolved by a sequence of blow-ups along nonsingular
centres \cite{hironaka-1964}. In positive characteristic this is known only in low dimension.
\end{remark}

\section{Exercises}\label{varieties:section-exercises}

\begin{exercise}\label{varieties:exercise-affine-line-punctured}
Show that $\Aa^1$ minus finitely many points is affine, and determine its coordinate ring.
\end{exercise}

\begin{solution}
If the points are $a_1, \ldots, a_m$, the set is $D(f)$ for $f = \prod_i(x - a_i)$, which is affine with coordinate ring $k[x]_f$ by
Example~\ref{varieties:example-morphisms}(1).
\end{solution}

\begin{exercise}\label{varieties:exercise-conics}
Let $\mathrm{char}\,k \neq 2$. Show that every irreducible conic $Z(q) \subseteq \PP^2$ ($q$ a quadratic form) is isomorphic to $\PP^1$.
\end{exercise}

\begin{solution}
Diagonalise $q$ by a linear change of coordinates. It has rank $3$, since a quadratic form of rank $\leq 2$ over an algebraically closed field
factors into linear forms. Over $k$ it is then equivalent to $xz - y^2$. The morphism $\PP^1 \to Z(xz - y^2)$, $[s : t] \mapsto [s^2 : st : t^2]$, is
the Veronese map $v_2$, an isomorphism onto its image (Example~\ref{varieties:example-projective}). Its inverse is given by $[x : y] $ on
$\{x \neq 0\}$ and by $[y : z]$ on $\{z \neq 0\}$.
\end{solution}

\begin{exercise}\label{varieties:exercise-cone-singular}
Let $\mathrm{char}\,k \neq 2$. Show that the cone $Z(x^2 + y^2 - z^2) \subseteq \Aa^3$ is singular only at the origin, and compute its tangent space
there.
\end{exercise}

\begin{solution}
The Jacobian is $(2x, 2y, -2z)$, which vanishes only at $0$, and $\dim = 2$, so the cone is nonsingular away from $0$
(Theorem~\ref{varieties:theorem-jacobian}). At $0$ the Jacobian is zero, so $T_0Y = k^3$ has dimension $3 > 2$.
\end{solution}

\begin{exercise}\label{varieties:exercise-hypersurfaces-meet}
Show that $n$ hypersurfaces in $\PP^n$ have a common point, and that the Segre quadric $Z(xw - yz) \subseteq \PP^3$ contains two families of lines
such that two lines of the same family are disjoint and two lines of different families meet in one point.
\end{exercise}

\begin{solution}
The first statement is Corollary~\ref{varieties:corollary-intersections-nonempty} with $Y = \PP^n$. The quadric is the image of the Segre map
$\PP^1 \times \PP^1 \to \PP^3$, $([a_0 : a_1], [b_0 : b_1]) \mapsto [a_0b_0 : a_0b_1 : a_1b_0 : a_1b_1]$, which is a bijection onto it. The images of
$\{a\} \times \PP^1$ and $\PP^1 \times \{b\}$ are lines. Two lines of the same family are images of disjoint sets, and $\{a\} \times \PP^1$ meets
$\PP^1 \times \{b\}$ exactly in $(a, b)$.
\end{solution}

\begin{exercise}\label{varieties:exercise-twisted-cubic-ideal}
Show that the projective twisted cubic $C = v_3(\PP^1) \subseteq \PP^3$ is cut out by the three quadrics $xz - y^2$, $yw - z^2$, $xw - yz$, but not by
any two of them.
\end{exercise}

\begin{solution}
The three quadrics say that the matrix $\left(\begin{smallmatrix}x & y & z \\ y & z & w\end{smallmatrix}\right)$ has rank $\leq 1$. Such a point with
$x \neq 0$ is $[1 : t : t^2 : t^3]$ with $t = y/x$, and with $x = 0$ it is $[0 : 0 : 0 : 1]$, so the common zeros are exactly $C$. The quadrics
$xz - y^2$ and $yw - z^2$ also vanish on the line $y = z = 0$, which is not contained in $C$. For the other pairs, $\{xz = y^2, xw = yz\}$ contains the
line $x = y = 0$, and $\{yw = z^2, xw = yz\}$ contains the line $z = w = 0$.
\end{solution}
